% ================================================================
% Các package cần cho phần này đã được khai báo trong preamble của report.tex:
%   \usepackage{algorithm}
%   \usepackage{algpseudocode}
%   \usepackage{booktabs}
% ================================================================

% ================================================================
\subsection{Nhận diện thực thể theo ngữ cảnh (Context Detector)}
% ================================================================

Thành phần nhận diện theo ngữ cảnh nhận đầu vào là một đoạn văn bản thô và
trả về danh sách các thực thể có dạng \texttt{type}, \texttt{text},
\texttt{start}, \texttt{end}. Khác với nhóm nhận diện bằng biểu thức chính quy,
nhóm này không chỉ dựa vào cấu trúc ký tự cố định mà còn khai thác các từ khóa,
vị trí xuất hiện, chữ hoa và quan hệ giữa các từ trong cùng câu.

Các loại thực thể được xử lý gồm tên người (\texttt{NAME}), tổ chức
(\texttt{ORGANIZATION}), địa điểm (\texttt{LOCATION}), tên người dùng
(\texttt{USERNAME}) và địa chỉ (\texttt{ADDRESS}). Về mặt thuật toán, quá trình
nhận diện được tổ chức theo luồng: chuẩn hóa ngữ cảnh đầu vào, sinh candidate,
chấm điểm hoặc kiểm tra mẫu, sau đó trả về các span thực thể không chồng lấn.

% ----------------------------------------------------------------
\subsubsection{Nhận diện tên người, tổ chức và địa điểm}
% ----------------------------------------------------------------

Đầu vào của thuật toán là văn bản đã được chia thành các câu cùng với vị trí của
từng câu trong văn bản gốc. Đầu ra là các span được gán một trong ba nhãn:
\texttt{NAME}, \texttt{ORGANIZATION} hoặc \texttt{LOCATION}. Ba loại thực thể
này đều có thể xuất hiện dưới dạng cụm danh từ riêng viết hoa, vì vậy hệ thống
không chạy ba bộ phát hiện tách rời. Thay vào đó, thuật toán sinh candidate một
lần, chấm điểm candidate đó theo cả ba nhãn, rồi quyết định nhãn cuối cùng ở
bước sau cùng.

Trước khi phân tích ngữ cảnh, văn bản đầu vào được tách thành các câu riêng
lẻ. Đây là bước quan trọng vì điểm ngữ cảnh (pronoun, nghề nghiệp, trigger)
chỉ có ý nghĩa trong phạm vi một câu.

Thách thức chính là phân biệt dấu chấm kết thúc câu với dấu chấm trong các
trường hợp khác. Thuật toán tách câu xử lý điều này theo quy tắc: dấu
\texttt{.} chỉ được coi là ranh giới câu khi \textit{đồng thời} thỏa ba điều
kiện:

\begin{enumerate}[nosep]
    \item Token kết thúc bằng \texttt{.} không thuộc danh sách viết tắt đã biết
          (\texttt{Dr.}, \texttt{Mr.}, \texttt{St.}, \texttt{Apt.}, \texttt{Inc.}, v.v.).
    \item Phần văn bản ngay sau \texttt{.} không bắt đầu bằng tên miền cấp cao
          (\texttt{com}, \texttt{edu}, \texttt{vn}, \texttt{org}, v.v.) --- nhằm
          tránh split nhầm địa chỉ email và URL.
    \item Sau dấu \texttt{.} là khoảng trắng hoặc hết chuỗi.
\end{enumerate}

Dấu \texttt{!} và \texttt{?} luôn được coi là kết thúc câu vô điều kiện.

\subsubsection*{Bước 1: xử lý trigger mạnh cho tên người}

Thuật toán trước hết tìm các cụm từ báo hiệu trực tiếp rằng phần đứng sau là
tên người. Các trigger này không tạo ra một luồng nhận diện riêng, mà được lưu
như một tín hiệu điểm mạnh cho nhãn \texttt{NAME}. Trigger được chia thành hai
nhóm:

\begin{itemize}[nosep]
    \item \textbf{NAME\_TRIGGERS}: cụm từ chỉ thẳng tên chủ thể, gồm
          \texttt{"my name is"}, \texttt{"full name"}, \texttt{"name:"}.
    \item \textbf{\_LABEL\_TRIGGERS}: nhãn hay đứng trước tên trong văn bản có
          cấu trúc, gồm khoảng 25 mẫu như \texttt{"contact information:"},
          \texttt{"author:"}, \texttt{"submitted by:"}, \texttt{"regards,"},
          \texttt{"sincerely,"}, v.v.
\end{itemize}

Khi tìm thấy trigger, thuật toán thu thập 1--4 token liên tiếp thỏa:
bắt đầu bằng chữ hoa, không chứa chữ số, không chứa ký tự đặc biệt (ngoại
trừ dấu gạch ngang \texttt{-} vì tên có thể có dạng \texttt{Mary-Jane}),
không phải stopword, và không mang đuôi động từ sau token đầu tiên
(\texttt{-ing}, \texttt{-tion}, \texttt{-ed}, v.v.). Span tìm được sẽ nhận điểm
\texttt{NAME} rất cao nếu trùng với candidate ở bước chấm điểm.

\subsubsection*{Bước 2: tìm candidate bằng linear search và rule-based filtering}

Với mỗi câu, thuật toán sinh các candidate danh từ riêng từ những cụm token viết
hoa. Candidate được dùng chung cho cả ba nhãn \texttt{NAME},
\texttt{ORGANIZATION} và \texttt{LOCATION}, thay vì viết lại ba hàm quét riêng.

Một candidate hợp lệ phải bắt đầu bằng chữ hoa, không chứa chữ số, không chứa
ký tự đặc biệt, không phải stopword, không phải đại từ, không phải từ mở đầu câu
thường gặp, không phải title prefix và không phải chức danh công việc. Thuật
toán cũng cho phép một số connector như \texttt{of}, \texttt{the},
\texttt{and} để bắt các cụm như \texttt{University of California}.

\subsubsection*{Bước 3: scoring candidate theo ba nhãn}

Chương trình sử dụng thuật toán Heuristic Rule-Based Scoring để gán điểm cho mỗi candidate theo ba nhãn. Mỗi nhãn có một tập hợp các quy tắc riêng. Ví dụ: tần suất xuất hiện của các từ đặc biệt, cấu trúc cố định,$\dots$

\subsubsection*{Output}
Sau khi có đủ ba điểm, thuật toán chọn đúng một nhãn cuối cùng cho candidate.
Nhãn có điểm cao nhất và đạt ngưỡng sẽ được giữ. Nếu nhiều nhãn bằng điểm,
thuật toán dùng thứ tự tie-break ổn định để tránh một span được xuất ra nhiều
lần với nhiều tag khác nhau. Vì vậy, một span như \texttt{OpenAI} trong câu
\texttt{I work at OpenAI.} được quyết định là \texttt{ORGANIZATION}, không đồng
thời xuất hiện thêm dưới dạng \texttt{NAME} hoặc \texttt{LOCATION}.

Sau khi mỗi candidate đã có nhãn cuối cùng, hệ thống loại trùng và overlap theo
priority nội bộ, điểm, vị trí bắt đầu và độ dài span trong phạm vi các kết quả
do context detector sinh ra. Kết quả này vẫn có thể được đưa sang merger tổng
thể để xử lý overlap với các thực thể từ regex detector.

\begin{algorithm}[H]
    \caption{Thuật toán nhận diện NAME, ORGANIZATION và LOCATION}
    \label{alg:detect_context_entities}
    \begin{algorithmic}[1]
        \Require Văn bản đầu vào $T$
        \Ensure Danh sách thực thể $E$

        \State $E \leftarrow \emptyset$
        \State $lower \leftarrow \text{lowercase}(T)$

        \Statex \hspace{-\algorithmicindent}\textbf{// Bước 1: Ghi nhận trigger mạnh cho NAME}
        \For{mỗi trigger $tr$ trong \textsc{NameTriggers} $\cup$ \textsc{LabelTriggers}}
        \For{mỗi vị trí $pos$ của $tr$ trong $lower$}
        \State Thu thập 1--4 token hoa liên tiếp sau $pos$ vào $tokens$
        \If{$tokens \neq \emptyset$}
        \State Lưu span $\text{join}(tokens)$ vào $strongNameSpans$ với điểm 100
        \EndIf
        \EndFor
        \EndFor
        \Statex \hspace{-\algorithmicindent}\textbf{// Bước 2--4: Candidate chung, scoring ba nhãn, chọn tag}
        \State $sentences \leftarrow \textsc{SplitSentences}(T)$
        \For{mỗi câu $s$ trong $sentences$}
        \For{mỗi candidate $c$ trong \textsc{ExtractEntityCandidates}($s$)}
        \State $nameScore \leftarrow \textsc{ScoreName}(c, s, strongNameSpans)$
        \State $orgScore \leftarrow \textsc{ScoreOrganization}(c, s)$
        \State $locScore \leftarrow \textsc{ScoreLocation}(c, s)$
        \State $tag \leftarrow \textsc{ChooseBestTag}(nameScore, orgScore, locScore)$
        \If{$tag$ tồn tại}
        \State Thêm $c$ vào $E$ với nhãn $tag$ và điểm tương ứng
        \EndIf
        \EndFor
        \EndFor

        \State \Return $\textsc{ResolveEntities}(E)$
    \end{algorithmic}
\end{algorithm}

% ----------------------------------------------------------------
\subsubsection{Nhận diện địa chỉ}
% ----------------------------------------------------------------

Đầu vào của thuật toán nhận diện địa chỉ là văn bản gốc và danh sách span email
đã được xác định trước. Đầu ra là các span \texttt{ADDRESS}. Địa chỉ trong tập
dữ liệu thường tuân theo cấu trúc gồm số nhà, tên đường, loại đường, và hậu tố
hướng hoặc đơn vị căn hộ tùy chọn. Thuật toán sử dụng hai cơ chế bổ sung cho
nhau.

\subsubsection*{Bước 1: nhận diện địa chỉ bằng Regex pattern}
Biểu thức chính quy \texttt{\_ADDRESS\_RE} mô tả cấu trúc địa chỉ theo bốn
thành phần nối tiếp:

\begin{enumerate}[nosep]
    \item \textbf{Số nhà}: \verb|\b(\d{1,5})(?!\d)\s+| ---
          từ 1 đến 5 chữ số, kèm negative lookahead \verb|(?!\d)| để đảm
          bảo đây không phải là một phần của số dài hơn (loại trừ số điện
          thoại, mã bưu điện dài).

    \item \textbf{Tên đường}: \verb|((?:[A-Z0-9][A-Za-z0-9]*\s+)+?)| ---
          một hoặc nhiều từ bắt đầu bằng chữ hoa hoặc chữ số, sử dụng
          non-greedy để nhường ưu tiên cho thành phần loại đường phía sau.

    \item \textbf{Loại đường}: danh sách khoảng 30 từ được sắp xếp từ dài
          đến ngắn để tránh match sớm, bao gồm dạng đầy đủ
          (\texttt{Boulevard}, \texttt{Terrace}, \texttt{Avenue},
          \texttt{Street}, \texttt{Drive}, \texttt{Road}, v.v.) và dạng viết
          tắt (\texttt{Blvd.}, \texttt{Ave.}, \texttt{Rd.}, \texttt{St.}, v.v.).

    \item \textbf{Hướng và đơn vị} (tùy chọn):
          hướng gồm tám giá trị \texttt{North}, \texttt{South}, \texttt{East},
          \texttt{West} và các tổ hợp; đơn vị gồm \texttt{Suite},
          \texttt{Apt.}, \texttt{Apartment}, \texttt{Unit} kèm số phòng.
\end{enumerate}

Trước khi kiểm tra kết quả regex, tất cả span email trong văn bản được xác định
trước; bất kỳ span địa chỉ nào nằm hoàn toàn bên trong span email sẽ bị bỏ qua
để tránh nhầm lẫn phần tên miền của địa chỉ email với địa chỉ nhà.

\subsubsection*{Bước 2: nhận diện địa chỉ bằng trigger keyword}
Để bổ sung cho regex chính, đặc biệt với các địa chỉ được viết theo dạng nhãn
hoặc clause theo trigger, module sử dụng danh sách
trigger bao gồm \texttt{street}, \texttt{road}, \texttt{avenue},
\texttt{district}, \texttt{city}, \texttt{st.}, \texttt{rd.}, \texttt{ave.}.
Khi tìm thấy trigger, thuật toán chỉ tiếp tục nếu ngay sau trigger có dấu
\texttt{:} hoặc cụm \texttt{is}. Sau đó hệ thống lấy clause đến dấu câu gần
nhất, loại bỏ noise words đầu clause (\texttt{at}, \texttt{is},
\texttt{the}, \texttt{located}, \texttt{live}, \texttt{reside}, v.v.), và kiểm
tra thêm: số nhà ở đầu clause không được vượt quá 5 chữ số, span không nằm
trong email, và span không được trùng lặp với kết quả từ cơ chế regex.

\begin{algorithm}[H]
    \caption{Thuật toán nhận diện địa chỉ}
    \label{alg:detect_address}
    \begin{algorithmic}[1]
        \Require Văn bản đầu vào $T$
        \Ensure Danh sách thực thể địa chỉ $E$

        \State $E \leftarrow \emptyset$
        \State $emailSpans \leftarrow \textsc{FindEmailSpans}(T)$

        \Statex \hspace{-\algorithmicindent}\textbf{// Cơ chế 1: Regex pattern}
        \For{mỗi match $m$ của \textsc{AddressRE} trong $T$}
        \If{$m$ nằm trong một email span} \textbf{continue} \EndIf
        \State Thêm span của $m$ vào $E$
        \EndFor

        \Statex \hspace{-\algorithmicindent}\textbf{// Cơ chế 2: Trigger keyword}
        \For{mỗi trigger $tr$ trong \textsc{AddressTriggers}}
        \For{mỗi vị trí $pos$ của $tr$ trong $T$}
        \If{sau $tr$ không phải \texttt{:} hoặc \texttt{is}} \textbf{continue} \EndIf
        \State $clause \leftarrow \textsc{ExtractClause}(T, pos)$
        \State $addr \leftarrow \textsc{StripLeadingNoise}(clause)$
        \If{số nhà đầu $addr$ có nhiều hơn 5 chữ số} \textbf{continue} \EndIf
        \If{$addr$ nằm trong email span} \textbf{continue} \EndIf
        \If{$addr$ overlap với span đã có trong $E$} \textbf{continue} \EndIf
        \State Thêm span của $addr$ vào $E$
        \EndFor
        \EndFor

        \State \Return $\textsc{Deduplicate}(E)$
    \end{algorithmic}
\end{algorithm}

% ================================================================
\subsection{Tổng hợp kết quả}
% ================================================================

Đầu vào của bước hợp nhất là nhiều danh sách thực thể được sinh ra từ các nhóm
nhận diện khác nhau. Đầu ra là một danh sách duy nhất, đã loại trùng và không
còn các span chồng lấn. Thuật toán hợp nhất là bước quyết định kết quả cuối
cùng trước khi văn bản được mã hoá.

\subsubsection{Vấn đề overlap}

Hai thực thể được coi là \textit{chồng lấn} (overlap) nếu khoảng ký tự của
chúng giao nhau, tức là $[s_1, e_1)$ và $[s_2, e_2)$ thỏa $s_1 < e_2$ và
$e_1 > s_2$. Overlap xảy ra trong nhiều tình huống thực tế:

\begin{itemize}[nosep]
    \item Địa chỉ email \texttt{kazuosun@hotmail.net} bị regex\_detector bắt
          đúng là \texttt{EMAIL}, trong khi các detector khác có thể tạo span
          chồng lấn trên một phần của chuỗi này.
    \item Một chuỗi như \texttt{123 Main Street} có thể được nhận là
          \texttt{ADDRESS}, đồng thời phần \texttt{Main Street} cũng có thể
          được nhận là \texttt{LOCATION}.
    \item Một cụm như \texttt{OpenAI Research Lab} có thể được chấm điểm theo
          nhiều nhãn ngữ cảnh; khi đó detector cần chọn một nhãn cuối cùng
          trước khi đưa kết quả sang bước hợp nhất.
\end{itemize}

\subsubsection{Hệ thống ưu tiên}

Để giải quyết overlap một cách nhất quán, merger định nghĩa một thứ tự ưu tiên
tuyến tính cho tám loại thực thể:

\begin{center}
    \texttt{EMAIL} $>$ \texttt{URL} $>$ \texttt{PHONE} $>$
    \texttt{ADDRESS} $>$ \texttt{ORGANIZATION} $>$ \texttt{LOCATION} $>$
    \texttt{NAME} $>$ \texttt{USERNAME}
\end{center}

Thứ tự này phản ánh độ tin cậy của từng detector: \texttt{EMAIL}, \texttt{URL}
và \texttt{PHONE} được phát hiện bằng regex có độ chính xác cao và cấu trúc
rõ ràng nên được ưu tiên hơn; \texttt{ADDRESS} thắng \texttt{LOCATION} để tránh
mất địa chỉ đầy đủ; \texttt{ORGANIZATION}, \texttt{LOCATION}, \texttt{NAME} và
\texttt{USERNAME} dựa nhiều hơn vào ngữ cảnh nên được xếp sau các entity có cấu
trúc rõ.

\subsubsection{Thuật toán Greedy Priority-First}

Thuật toán hợp nhất thực hiện giải quyết overlap theo chiến lược tham lam ưu
tiên (greedy priority-first):

\begin{enumerate}[nosep]
    \item Gộp tất cả thực thể từ mọi detector thành một danh sách duy nhất.
    \item Sắp xếp danh sách theo ba tiêu chí lần lượt: (a)~thứ hạng ưu tiên
          tăng dần (loại ưu tiên cao xếp trước), (b)~chỉ số \texttt{start}
          tăng dần, (c)~độ dài span giảm dần (span dài hơn xếp trước khi cùng
          điểm xuất phát).
    \item Duyệt tuần tự qua danh sách đã sắp xếp. Với mỗi thực thể, giữ lại
          nếu span của nó không chồng lấn với bất kỳ thực thể đã giữ nào;
          bỏ qua nếu ngược lại.
    \item Sắp xếp lại danh sách kết quả theo \texttt{start} để đảm bảo thứ tự
          xuất hiện trong văn bản.
\end{enumerate}

\begin{algorithm}[H]
    \caption{Thuật toán hợp nhất và giải quyết overlap (Greedy Priority-First)}
    \label{alg:merger}
    \begin{algorithmic}[1]
        \Require Các nhóm thực thể $G_1, G_2, \ldots, G_k$ từ các detector
        \Ensure Danh sách thực thể không overlap, sắp theo vị trí

        \State $all \leftarrow G_1 \cup G_2 \cup \cdots \cup G_k$

        \State Sắp xếp $all$ theo $\bigl(\text{priority}(type),\ start,\ -(end - start)\bigr)$ tăng dần

        \State $kept \leftarrow \emptyset$
        \State $occupied \leftarrow \emptyset$ \quad \textit{// Tập các span $(s, e)$ đã giữ}

        \For{mỗi thực thể $ent$ trong $all$}
        \State $s \leftarrow ent.start$,\quad $e \leftarrow ent.end$
        \If{$\nexists\, (os, oe) \in occupied$ sao cho $s < oe$ và $e > os$}
        \State $kept \leftarrow kept \cup \{ent\}$
        \State $occupied \leftarrow occupied \cup \{(s, e)\}$
        \EndIf
        \EndFor

        \State Sắp xếp $kept$ theo $start$ tăng dần
        \State \Return $kept$
    \end{algorithmic}
\end{algorithm}

\subsubsection{Minh họa trường hợp overlap}

Xét đoạn văn bản sau:

\begin{center}
    \textit{``I live at 123 Main Street.''}
\end{center}

Trong đoạn này, \texttt{ADDRESS} và \texttt{LOCATION} có thể cùng nhận diện
vùng văn bản liên quan đến \texttt{Main Street}. Danh sách trước khi hợp nhất
có dạng như Bảng~\ref{tab:overlap_before}.

\begin{table}[H]
    \centering
    \caption{Danh sách thực thể trước khi hợp nhất}
    \label{tab:overlap_before}
    \renewcommand{\arraystretch}{1.3}
    \begin{tabular}{l l l c}
        \toprule
        \textbf{Detector} & \textbf{Loại}     & \textbf{Text}            & \textbf{Ưu tiên} \\
        \midrule
        context\_detector & \texttt{ADDRESS}  & \texttt{123 Main Street} & 4                \\
        context\_detector & \texttt{LOCATION} & \texttt{Main Street}     & 6                \\
        \bottomrule
    \end{tabular}
\end{table}

Sau khi sắp xếp theo ưu tiên và áp dụng thuật toán greedy, merger giữ
\texttt{123 Main Street} vì \texttt{ADDRESS} có priority cao hơn
\texttt{LOCATION}. Candidate \texttt{Main Street} bị loại do overlap với span
địa chỉ đầy đủ đã được giữ.
Kết quả cuối cùng được trình bày trong Bảng~\ref{tab:overlap_after}.

\begin{table}[H]
    \centering
    \caption{Danh sách thực thể sau khi hợp nhất}
    \label{tab:overlap_after}
    \renewcommand{\arraystretch}{1.3}
    \begin{tabular}{l l l}
        \toprule
        \textbf{Loại}    & \textbf{Text}            & \textbf{Vị trí}  \\
        \midrule
        \texttt{ADDRESS} & \texttt{123 Main Street} & start=10, end=25 \\
        \bottomrule
    \end{tabular}
\end{table}

Kết quả này sau đó được đưa vào module \texttt{anonymizer} để thay thế các
span PII bằng placeholder tương ứng, tạo ra văn bản đã được ẩn danh hóa.
